\chapter{Related Work}
\label{sec:related-work}

\section{Background}
Research on free-rider (FR) behavior in Federated Learning (FL) can be organized into five broad categories: (1) foundational work that formalizes free-rider attacks and analyzes their effect on convergence, (2) increasingly stealthy attack constructions that exploit the geometric and statistical properties of honest updates, (3) server-side detection and audit-based mitigation mechanisms, (4) detection methods that repurpose privacy-attack techniques, and (5) hardware-based and cryptographic verification schemes. This section reviews each category in turn, compares the underlying assumptions and effectiveness of representative methods, and concludes by identifying the gap that motivates the approach proposed in this thesis.

\section{Foundational Characterization of Free-Rider Attacks}
\label{subsec:foundational}

Fraboni et al. \cite{fraboni2021freerider} provided the first theoretical and experimental analysis of free-rider attacks against iterative parameter-aggregation schemes such as FedAvg and FedProx. Modeling stochastic gradient descent as a continuous-time process, the authors formally proved that a free rider's fabricated update can be constructed to provably converge toward the aggregated model produced by fair participants, meaning the attacker obtains a fully-trained global model without contributing any data. They first showed that the most naive attack—simply returning the previous global model unchanged—is trivially detectable by basic server-side checks (e.g., zero update norm). To evade such checks, they proposed a \emph{disguised free-rider} strategy that adds time-decaying Gaussian noise to the previous global model, with the noise variance calibrated to the empirical standard deviation of the global gradient so that the fabricated update statistically mimics the natural shrinkage of SGD updates as training converges. Their experiments across both IID and non-IID settings showed that this disguised attack is substantially harder to detect than the naive variant, and they concluded with practical recommendations for FL deployments in sensitive, high-value domains such as healthcare.

Complementing this, Lin et al. \cite{lin2019freeriders} were the first to formally introduce the notion of free-rider attacks in FL and to jointly study attack construction and detection. They proposed the \emph{delta weights} attack, in which a free rider constructs its fake update by adding calibrated noise to the previous round's global gradient rather than to the model weights directly, producing updates that more closely resemble genuine gradient-based contributions than naive random-weight submissions. On the defense side, they proposed \textbf{STD-DAGMM}, a detection method that augments a Deep Autoencoding Gaussian Mixture Model (DAGMM) with the standard deviation of each client's update as an additional feature, framing free-rider detection as a high-dimensional anomaly-detection problem. They further extended their analysis to scenarios involving multiple simultaneous free riders and the presence of differential privacy noise, showing that the interaction between privacy-preserving noise and free-rider disguise strategies further complicates detection—an open problem they explicitly flagged for future research.

Together, these two works established that (i) free-riding is a concretely realizable and theoretically grounded threat rather than a hypothetical one, and (ii) naive server-side aggregation, and even simple statistical checks, are insufficient once an attacker calibrates its fabricated update to mimic the stochastic behavior of genuine local training.

\section{Advanced and Stealthy Free-Rider Attacks}
\label{subsec:advanced-attacks}

Building directly on the disguising strategies above, Zhu et al. \cite{zhu2023advanced} introduced a more sophisticated \emph{advanced free-rider attack} that explicitly targets the geometric properties exploited by feature-based defenses. The authors analyzed the expected cosine similarity and $\ell_2$-norm relationships between honest client updates and the converging global model in IID settings, deriving closed-form approximations for how these quantities evolve as training progresses. Using these derived relationships, their attack injects calibrated Gaussian noise into only a \emph{subset} of the global gradient's parameters, with the noise magnitude explicitly scaled so that the resulting fake update reproduces the cosine-similarity and norm statistics expected of a genuine client update at that point in training. Because this attack is constructed to match the exact geometric signature that earlier detection mechanisms (including cosine-similarity and $\ell_2$/standard-deviation-based checks) rely on, it was shown—both in the original work and in later independent evaluations \cite{recasens2024frida}—to substantially degrade the effectiveness of feature-based detectors, particularly in IID scenarios where honest client updates are naturally similar to one another and thus offer the attacker less statistical cover to hide \emph{behind} but also less distinctiveness to be caught \emph{by}.

This line of work demonstrates a critical limitation of static, feature-based defenses: any detection rule based on a fixed set of update statistics (norm, standard deviation, cosine similarity) can, in principle, be reverse-engineered and matched by an adaptive attacker with knowledge of the aggregation dynamics. It motivates the need for detection mechanisms that target properties of genuine training which are fundamentally difficult to fabricate, rather than statistical proxies that a sufficiently informed adversary can approximate.

\section{Server-Side Detection and Audit-Based Mitigation}
\label{subsec:server-side}

A second line of work designs server-side detection pipelines that combine anomaly detection with auxiliary contribution or reputation signals. Wang et al. proposed \textbf{FRAD} \cite{wang2024frad}, a free-rider attack detection mechanism targeted at FL deployments in the Artificial Intelligence of Things (AIoT), where device heterogeneity is pronounced. FRAD combines a DAGMM-based anomaly detector (extending the STD-DAGMM approach of Lin et al.) with an explicit \emph{contribution model} that estimates each device's expected contribution from its computing resources, communication cost, and data quality, and a \emph{reputation model} based on the PageRank algorithm that iteratively re-weights clients according to their historical behavior. By fusing statistical anomaly detection with resource- and reputation-aware contribution modeling, FRAD is explicitly designed to select benign nodes for participation even under conditions of information asymmetry between the server and heterogeneous IoT clients. However, because its detection backbone remains a DAGMM-style autoencoder, it inherits the same scalability limitation identified in later work \cite{recasens2024frida}: the memory and computational cost of training an autoencoder over every client's full (flattened) update grows with both model size and the number of clients, which becomes prohibitive for larger neural networks or large client populations.

Wang et al. also proposed \textbf{PASS} \cite{wang2023pass}, a Parameter Audit-based Secure and fair FL Scheme, motivated by the observation that prior defenses generally assumed adversaries constitute a minority of clients and focused narrowly on \emph{anonymous} free riders (clients lacking data or compute) while losing effectiveness against \emph{selfish} free riders (clients that possess adequate resources but strategically withhold them). PASS combines a privacy-preserving auditing strategy—injecting weak differential-privacy noise via a Gaussian mechanism together with a parameter-pruning step to protect client data during the audit—with a novel contribution-evaluation method used to assess each client's genuine performance. The scheme's experimental evaluation demonstrated that PASS achieves privacy protection comparable to state-of-the-art baselines with limited accuracy loss, attains a higher defense success rate, lower false-positive rate, and higher F1-score than prior methods against both anonymous and selfish free-rider attacks, and—notably—remains effective even when adversarial clients constitute a \emph{majority} of the population (F1-scores of 89\% and 87\% against anonymous and selfish attacks, respectively, even when adversaries exceed 50\% of clients). This robustness to majority-adversary settings distinguishes PASS from earlier defenses, though its parameter-auditing procedure introduces non-trivial computational overhead at the server compared to lightweight statistical baselines.

Collectively, FRAD and PASS demonstrate that meaningful detection is achievable using only server-observable signals augmented with lightweight contribution modeling, aligning with FL's privacy-preserving ethos. However, both rely on carefully engineered statistical or audit pipelines whose effectiveness depends on assumptions (e.g., DAGMM's density-estimation quality, or the calibration of PASS's audit noise) that may require retuning as data distributions, model architectures, and free-rider strategies evolve.

\section{Privacy-Attack-Based Detection}
\label{subsec:privacy-attack-based}

A more recent line of work, exemplified by Recasens et al.'s \textbf{FRIDA} (Free-Rider detection using privacy Attacks) \cite{recasens2024frida}, departs from indirect statistical proxies and instead repurposes \emph{privacy-attack} methodologies to directly verify whether a client's update reflects genuine training. The key insight is that prior detection families—feature-based statistics (norm, standard deviation, cosine similarity), DAGMM-based anomaly detection (including FRAD), log-based auditing, and incentive/reputation mechanisms—all rely on \emph{indirect effects} of free-riding rather than on the root cause: the absence of genuine training on real client data. FRIDA instead proposes four detection mechanisms built on two families of privacy attack. The first family adapts \emph{membership inference attacks} (MIA), originally studied as a way to infer whether a record was part of a model's training data \cite{shokri2017membership}: the server distributes a small shared ``canary'' dataset that honest clients are required to additionally train on each round; a \emph{loss-based} detector then flags clients whose canary-set loss fails to decrease as expected, while a \emph{cosine-based} detector compares the cosine similarity of a client's gradient with canary-sample gradients against non-canary samples via a statistical test, flagging clients for whom the two distributions are indistinguishable (indicating no canary training occurred). The second family adapts \emph{property inference attacks} (PIA), which exploit information leakage from collaborative model updates to infer aggregate properties of a participant's local data \cite{melis2019exploiting}, to estimate each client's local label distribution from its submitted gradients without requiring any shared dataset or additional client-side training; a \emph{consistency-based} detector flags clients whose inferred label distribution changes anomalously across rounds, while a \emph{diversity-based} detector flags clients whose inferred distribution is best explained as a combination of pure noise and the global label distribution—the signature of a fabricated, rather than genuinely trained, update.

The authors' extensive evaluation across CIFAR-10 and CIFAR-100 with AlexNet and VGG-19 architectures, under both IID and non-IID data (Dirichlet-distributed with $\alpha \in \{0.1, 1.0\}$), and against all three attack strategies discussed above (Fraboni et al., Lin et al., and Zhu et al.), showed that FRIDA matches or outperforms feature-based and DAGMM-based baselines in nearly every setting, with especially pronounced advantages in IID scenarios and on more complex datasets—precisely the regimes in which the advanced attack of Zhu et al. is most effective at evading feature-based detectors. Notably, PIA-based detection was shown to remain robust even when honest clients apply local differential privacy to their updates, since it exploits dataset-level statistical properties rather than fine-grained individual-level signals, whereas MIA-based detection degrades sharply under such privacy noise. However, FRIDA's design assumes the server has visibility into individual client updates, which is fundamentally in tension with secure-aggregation-based privacy protection; the authors propose a two-server architecture as a partial mitigation but acknowledge this adds system complexity, and the cosine-based MIA variant was shown to be evadable by an adaptive, resource-endowed free rider willing to train specifically on the canary set. The MIA-based variants also introduce non-trivial per-round computational cost relative to lightweight feature-based statistics, and PIA-based detection shifts additional computational burden onto the server.

\section{Hardware-Based and Cryptographic Solutions}
\label{subsec:hardware-based}

In parallel with statistical and inference-based defenses, several works have explored hardware-based and cryptographic approaches. Trusted Execution Environments (TEEs), such as Intel SGX, have been proposed to provide attestable, tamper-evident proof that local training genuinely executed on a client device. Similarly, blockchain-based FL frameworks have incorporated proof-of-work or proof-of-training consensus mechanisms that require clients to demonstrate verifiable computational effort before their updates are accepted into aggregation, and secure-aggregation or zero-knowledge-proof-based schemes have been proposed to protect or verify client-side computation without exposing raw client data or full individual updates \cite{bonawitz2017practical}.

While such approaches can offer strong, cryptographically verifiable guarantees against fabricated updates, they carry substantial practical costs. TEE-based solutions require specialized hardware support that is unavailable across much of the heterogeneous device landscape typical of real-world FL (older mobile devices, low-cost IoT sensors), limiting deployability. Proof-of-work mechanisms impose significant computational and energy overhead that is particularly problematic for battery-constrained edge devices and works directly against FL's goal of enabling participation from resource-limited clients. Zero-knowledge-proof schemes, while avoiding the hardware-dependency problem, typically introduce substantial proof-generation and verification overhead that scales poorly with model size. None of these approaches inherently addresses the non-IID, dynamic nature of real-world client populations, and all introduce deployment complexity that has, to date, limited their adoption outside specialized or high-value settings.

\section{Other Notable Approaches}
\label{subsec:other-approaches}

Beyond the categories above, several related lines of work address free riding indirectly through contribution evaluation and incentive design. As noted by Recasens et al. \cite{recasens2024frida}, free-rider detection can be framed as an extreme case of the broader contribution-evaluation problem, for which the Shapley value is the game-theoretic gold standard; however, exact computation is combinatorially intractable, and approximations such as Data Shapley and task-specific Shapley valuation are generally designed to \emph{rank} contributors rather than precisely identify zero-contribution participants \cite{ghorbani2019data,jia2019efficient}, limiting their direct applicability to free-rider detection. Log-based auditing schemes detect free riders by inspecting client activity or stake logs rather than model updates directly, while incentive- and reputation-based frameworks allocate rewards or model access proportionally to accumulated trust across rounds. These approaches are complementary to update-based detection but operate on information beyond the raw model updates (e.g., detailed activity logs or reward accounting) and remain vulnerable to strategic clients who behave honestly during an initial trust-building phase before defecting once sufficient reputation has accrued.

\section{Summary and Research Gap}
\label{subsec:research-gap}

Table~\ref{tab:related-work-summary} summarizes the reviewed approaches according to their privacy posture, computational cost, and robustness to adaptive, stealthy free-rider strategies such as that of Zhu et al. \cite{zhu2023advanced}.

\begin{table}[htbp]
\centering
\caption{Comparative summary of existing free-rider mitigation approaches}
\label{tab:related-work-summary}
\begin{tabular}{p{3cm}p{3.6cm}p{2.1cm}p{2.1cm}p{2.4cm}}
\toprule
\textbf{Category} & \textbf{Representative Work} & \textbf{Privacy-Preserving} & \textbf{Lightweight} & \textbf{Robust to Advanced Attacks} \\
\midrule
Foundational \& statistical detection & Fraboni et al. \cite{fraboni2021freerider}; Lin et al. \cite{lin2019freeriders} (STD-DAGMM) & Yes & Moderate & Limited \\
Advanced attack construction & Zhu et al. \cite{zhu2023advanced} & N/A (attack) & N/A & N/A \\
Server-side audit-based & FRAD \cite{wang2024frad}; PASS \cite{wang2023pass} & Yes & Moderate & Limited--Moderate \\
Privacy-attack-based & FRIDA \cite{recasens2024frida} & Partial\footnotemark & No (MIA); Moderate (PIA) & Yes (esp. PIA) \\
Hardware/ cryptographic & TEE, PoW, zero-knowledge proofs & Yes & No & Yes \\
Contribution/ incentive-based & Shapley approximations; reputation systems & Yes & No & Partial \\
\bottomrule
\end{tabular}
\end{table}
\footnotetext{FRIDA requires server-side inspection of individual client updates, in tension with secure aggregation; the authors propose a two-server architecture as a partial mitigation, and show PIA-based detection remains effective under local differential privacy.}

Several patterns emerge from this review. First, purely statistical and feature-based methods \cite{fraboni2021freerider,lin2019freeriders} are lightweight and privacy-preserving but were explicitly shown to be circumventable by attacks constructed to match their target statistics \cite{zhu2023advanced}, and this vulnerability propagates into the audit-based methods (FRAD, PASS) that build on similar anomaly-detection backbones. Second, the privacy-attack-based approach of FRIDA \cite{recasens2024frida} achieves markedly better robustness against exactly this class of advanced, geometry-aware attacks, but its MIA variants require client compliance in training on server-issued canary data (an assumption that a resource-constrained or adversarially adaptive free rider may violate or exploit), and both MIA and PIA variants require the server to inspect individual, non-aggregated client updates—placing FRIDA in direct tension with secure aggregation, one of FL's core privacy-preserving mechanisms. Third, hardware- and cryptography-based methods offer strong guarantees largely independent of the free rider's statistical sophistication, but their cost, hardware dependency, and poor fit for heterogeneous, resource-constrained deployments make them impractical for many real-world FL scenarios. Finally, contribution- and incentive-based methods provide a longitudinal, economically grounded perspective but are either computationally intractable in exact form or vulnerable to strategic, time-varying adversarial behavior.

Across all categories, no reviewed method simultaneously satisfies three properties that are jointly necessary for practical deployment in large-scale, real-world FL systems: (i) \emph{privacy preservation}, in the sense of not requiring the server to inspect raw, individually attributable client updates or client data, and remaining compatible with mechanisms such as secure aggregation; (ii) \emph{lightweight operation}, avoiding the training of auxiliary models (autoencoders, shadow models) or the imposition of additional client-side training and communication burdens that scale poorly with model size or client population; and (iii) \emph{robustness to dynamic, adaptive adversaries}, remaining effective against free-rider strategies explicitly designed—as in Zhu et al. \cite{zhu2023advanced}—to evade previously published defenses, and against client populations and data distributions that shift over the course of training. Methods that are lightweight and privacy-preserving (feature-based statistics) sacrifice robustness to advanced attacks; methods that are robust and reasonably privacy-conscious (FRIDA) sacrifice lightweight operation and full compatibility with secure aggregation; and methods that are robust and lightweight in the security sense (hardware/cryptographic) sacrifice deployability, cost-effectiveness, and applicability to heterogeneous edge and IoT clients.

This unresolved three-way trade-off constitutes the central research gap motivating this thesis. The remainder of this work develops a free-rider detection and mitigation framework explicitly designed to occupy the previously unaddressed region of this design space: relying only on information that is either already available to the server or obtainable without violating standard FL privacy guarantees, avoiding the training of auxiliary detection models or reliance on client-side compliance with server-issued probing data, and evaluated explicitly against adaptive, geometry-aware free-rider strategies across dynamic, non-IID, and large-scale client populations.
